\TeoremaBox{Caracterización del Producto Semidirecto Interno}{
Sea $G$ un grupo. Sean $K \triangleleft G$ y $H \leq G$ subgrupos tales que $G = HK$ y $H \cap K = \{e\}$. Entonces
\begin{equation*}
G \cong H \ltimes_{\varphi} K
\end{equation*}
donde $\varphi$ es la acción por conjugación (inversa).
}

\begin{proof}
Definamos la función $f: H \ltimes_{\varphi} K \to G$ mediante la regla de correspondencia:
\begin{equation*}
f(a,b) = ab
\end{equation*}
para todo $(a,b) \in H \ltimes_{\varphi} K$.

Primero verifiquemos que $f$ es un homomorfismo.
Sean $(a,b), (a',b') \in H \ltimes_{\varphi} K$. Evaluamos $f$ en el producto de estos elementos:
\begin{align*}
f((a,b)(a',b')) &= f(aa', b^{a'}b') \\
&= (aa')(b^{a'}b')
\end{align*}
Recordemos que $b^{a'}$ denota la acción de $a'$ sobre $b$, que definimos como la conjugación inversa. Sustituyendo esto:
\begin{align*}
(aa')(b^{a'}b') &= a a' ((a')^{-1} b a') b' \\
&= a (a' (a')^{-1}) b a' b' \\
&= a e b a' b' \\
&= a b a' b'
\end{align*}
Por otro lado, calculamos el producto de las imágenes en $G$:
\begin{equation*}
f(a,b)f(a',b') = (ab)(a'b') = aba'b'
\end{equation*}
Dado que ambos desarrollos coinciden, concluimos que
\begin{equation*}
f((a,b)(a',b')) = f(a,b)f(a',b')
\end{equation*}

Ahora veamos que $f$ es inyectiva.
Supongamos que $(a,b) \in \ker(f)$. Entonces:
\begin{equation*}
f(a,b) = e \iff ab = e
\end{equation*}
De la igualdad $ab=e$, despejando $b$, obtenemos $b^{-1} = a$.
Sabemos que $a \in H$ y $b \in K$ (lo que implica $b^{-1} \in K$).
Por lo tanto:
\begin{equation*}
b^{-1} = a \in H \cap K
\end{equation*}
Por hipótesis, la intersección es trivial, es decir $H \cap K = \{e\}$. Esto implica que $a = e$.
Sustituyendo en la ecuación original, $b = a^{-1} = e$.
Así, $(a,b) = (e,e)$, por lo que el núcleo es trivial y $f$ es inyectiva.

Finalmente, veamos que $f$ es sobreyectiva.
Sea $g \in G$. Como $G = HK$, existen $a \in H$ y $b \in K$ tales que $g = ab$.
Luego:
\begin{equation*}
f(a,b) = ab = g
\end{equation*}
lo cual muestra que $f$ es sobreyectiva.

Resumiendo la construcción:
Dada una acción $\varphi: H \to \text{Aut}(K)$, formamos $(G, \cdot)$ con el producto
\begin{equation*}
(a,b)(a',b') = (aa', b^{a'}b')
\end{equation*}
Esto implica que $G = H \ltimes_{\varphi} K$ es un grupo donde identificamos $H \leq G$ y $K \triangleleft G$ tales que $G = HK$ y $H \cap K = \{e\}$.

Recíprocamente, si tenemos
\begin{equation*}
H \leq G, \quad K \triangleleft G, \quad G = HK
\end{equation*}
y
\begin{equation*}
H \cap K = \{e\}
\end{equation*}
entonces
\begin{equation*}
G \cong H \ltimes_{\varphi} K
\end{equation*}
\end{proof}

\begin{eje}
Sea $G$ un grupo de orden 8. Supongamos que $G$ no es abeliano.

Entonces $G$ tiene un elemento de orden 4.
En efecto, supongamos lo contrario (que $G$ no tiene elementos de orden 4). Entonces todos sus elementos distintos de la identidad tendrían orden 2. Es decir, $a^2 = e$ para todo $a \in G$.
Esto implicaría que $G$ es abeliano, pues para cualesquiera $a, b \in G$:
\begin{align*}
ab &= (ab)^{-1} \quad (\text{pues } (ab)^2=e) \\
&= b^{-1}a^{-1} \\
&= ba \quad (\text{pues } b^2=e, a^2=e \implies b=b^{-1}, a=a^{-1})
\end{align*}
Esto es una contradicción ($G$ no es abeliano). Por lo tanto, existe un elemento de orden 4.

Por lo tanto, $G$ tiene al menos un elemento de orden 4.
Sea $a \in G$ tal que $o(a)=4$. Definamos $K = \langle a \rangle$.
Como $|K|=4$ y $|G|=8$, el índice es $[G:K]=2$, lo que implica que $K \triangleleft G$.

Ahora analizamos los elementos que están fuera de $K$. Distinguimos dos casos:

Caso 1. Existe $b \in G \setminus K$ tal que $b^2 = e$.
Tomando $H = \langle b \rangle$, se tiene un subgrupo de orden 2.
Verificamos las condiciones:
\begin{itemize}
    \item $H \cap K = \{e\}$ (pues $b$ tiene orden 2 y el único elemento de orden 2 en $K$ es $a^2$, y $b \neq a^2$).
    \item $K \triangleleft G$.
    \item $G = HK$ (pues $|HK| = \frac{|H||K|}{|H \cap K|} = 8$).
\end{itemize}
Por el teorema anterior, tenemos que:
\begin{equation*}
G \cong H \ltimes_{\varphi} K
\end{equation*}

Para determinar el grupo, analizamos la acción $\varphi: H \to \text{Aut}(K)$.
Sabemos que $\text{Aut}(K) \cong (\mathbb{Z}_4)^\ast \cong \mathbb{Z}_2$, generado por la inversión $\theta(k) = k^{-1}$.
Si $\varphi$ fuera trivial, $G$ sería abeliano (contradicción).
Por lo tanto, $\varphi$ debe ser el isomorfismo que envía el generador de $H$ a la inversión en $K$.
Esta estructura corresponde al grupo Diédrico:
\begin{equation*}
G \cong D_4
\end{equation*}

Caso 2. Todo elemento de orden 2 está en $K$.
Supongamos que no existe ningún elemento de orden 2 fuera de $K$. Entonces, cualquier $b \in G \setminus K$ debe tener orden 4.
Si tomamos $H = \langle b \rangle$, entonces $H$ es un subgrupo de orden 4.
Observamos que $a^2$ es el único elemento de orden 2 en $K$, y $b^2$ es el único elemento de orden 2 en $H$. Como por hipótesis todos los elementos de orden 2 están en $K$, necesariamente $b^2 = a^2$.
Esto implica que $H \cap K \neq \{e\}$ (la intersección es al menos $\{e, a^2\}$).
Por lo tanto, en este caso $G$ no se puede descomponer como un producto semidirecto interno de un subgrupo de orden 4 y uno de orden 2 con intersección trivial.
(Nota: Este caso corresponde al grupo de los Cuaterniones, $Q_8$).

\end{eje}

